\section{Introduction}\label{sec:intro}

Stephen Hawking's \emph{A Brief History of Time} ends with a hope. If physics found a complete
theory and we could then understand why the universe exists, ``it would be the ultimate triumph
of human reason---for then we would know the mind of God'' \citep[ch.~11]{hawking1988}.
\emph{Algorithmic Agnosticism} argued that no account of everything can be known to be final, and that what
selects among the accounts the evidence leaves open is a choice \citep[§§3, 9]{taylor2026aa}.

This paper asks the parallel question of religious belief. It does not ask whether God exists. It
asks whether any worldview---theism, naturalism, the skeptical hypothesis, or the teaching of any
religious tradition---can be established as the true account of what is.

Its position has three parts. A worldview is a claim about what ultimately is, and an account
consistent with everything anyone has experienced makes it false. To establish a worldview would be
to show that the claim holds in every account the evidence leaves open. No one can do that, at any
time, and what selects among the worldviews the evidence leaves open is a choice. Stated for the worldviews at issue:

\begin{quote}
No argument from experience or reason establishes that any worldview, whether theism,
naturalism or the skeptical hypothesis, is the true account of what is. Any of them may be true,
but no one can establish that it is.
\end{quote}

The claim is one of parity. It is not that God does not exist, or that nature is all there is.

The parity matters because it cuts both ways. It removes demonstration from the theist and the
atheist alike. The proofs of natural theology and the disproofs offered against them fall under the
same verdict, so neither side can claim that reason has decided the question in its favor. It
leaves belief and faith untouched, and with them the question of what it is reasonable to believe;
the paper is silent on these. What the parity rules out is any argument's claim to settle the question for
everyone. An argument may persuade those who accept its premises and its standard of preference.
It does not settle the question for anyone who shares the same evidence and declines them.

The core claim has neighbors. \citet{oppy2006} argues that no argument about God is successful, by a
dialectical standard of success. \citet{mizrahi2017} argues that God's existence can be
established neither \emph{a priori} nor \emph{a posteriori}, a position he calls skeptheism.
\citet{alston1986} argues that the reliability of a doxastic practice cannot be shown without
epistemic circularity. \citet{hick1989} and \citet{mckim2001} argue that the universe is
religiously ambiguous. This paper differs in three ways. It derives its conclusions from the
premise and results of \emph{Algorithmic Agnosticism}, with a premise that extends its first result to worldviews. Its classification locates the
decisive element of the classic arguments in \cref{sec:arguments}. And its parity includes
naturalism and the skeptical hypothesis.

\Cref{sec:aa} restates the results of \emph{Algorithmic Agnosticism}. \Cref{sec:classification} gives the
definitions and the premises. \Cref{sec:arguments} takes up the classic arguments.
\Cref{sec:result} derives the conclusions. \Cref{sec:objections} answers objections, and
\cref{sec:conclusion} concludes.
